\section{Discussion}

The fixed-fold benchmark indicates that similarity-neighbour information can add predictive signal beyond the training-fold structural encoder, but only when its construction is constrained to the training partition. The full S1/S2 model had higher mean AUPR than the structural-only variant in 11/12 conditions, and 11 descriptive paired contrasts had a Benjamini--Hochberg-adjusted $p<0.05$. This pattern is consistent with the interpretation that local neighbourhood support and drug--protein consistency can complement graph-derived pair representations.

The shuffled control provides a narrower interpretation of this association. The full model had higher mean AUPR than the shuffled S1/S2 control in 12/12 conditions, including 11 descriptive paired contrasts with FDR-adjusted $p<0.05$. Because shuffling preserves the marginal side-feature values while breaking their alignment with each candidate pair, the comparison suggests that aligned side information, rather than only feature scale or dimensionality, contributes to the observed differences. It does not by itself establish a biological mechanism for any predicted interaction.

The result is not a claim of uniform superiority. The full model was not the highest-AUPR principal baseline in BioSNAP Cold protein, Human Cold protein, Human Cold pair, BindingDB Cold protein, BindingDB Cold pair. These settings should be treated as part of the method boundary: a similarity-neighbour summary may be less informative when representation coverage, dataset composition or the held-out entity structure differs from the conditions in which it is advantageous.

Several limitations follow directly from the evaluation design. The labels are curated DTIs, and the capped samples of unannotated pairs may include unknown positives; the reported metrics therefore describe the specified benchmark sampling distributions rather than the prevalence of interactions in a complete chemical--proteomic space. The study uses fixed molecular and protein feature caches, rather than learning representations from raw structures or sequences within each fold. Finally, the benchmark evaluates ranking and classification against annotated labels; experimental binding assays and prospective candidate selection are required before a prediction can support a pharmacological claim.
